\section{Step 6 Decode the remaining fifteen words}\label{sha-s06}

Advance both pointers by four bytes and decrement the remaining-word count. The branch returns to schedule\_decode while another source group remains. After the sixteenth store, \texttt{t0} points immediately beyond those sixteen words; load forty-eight as the number of expanded entries still to compute. The next region begins with exactly the predecessor words expansion needs.

\subsection{Write and inspect}
Replace only the temporary instructions between \texttt{\# BEGIN sch\_walk} and \texttt{\# END sch\_walk}. Keep the supplied labels outside that region.

\Needspace{10\baselineskip}
\MARSCodeLabel{MIPS assembly}
\begin{minted}[samepage=true]{mars-idea}
    daddiu $a0, $a0, 4
    daddiu $t0, $t0, 4
    addiu $t1, $t1, -1
    bne $t1, $zero, schedule_decode
    nop
    li $t1, 48
\end{minted}

Use recipe \texttt{s06-1} with the cumulative program through this edit. The exact recipe names the entry, prepared input and checkpoint. Observe the checkpoint before continuing to its temporary or normal return.

\subsection{Expected observation}
The first sixteen schedule words match the block; \texttt{t0} is 64 bytes beyond the supplied schedule address.

\Needspace{9\baselineskip}
\subsection{Check your understanding}
Why does the source advance by four?

\begin{itemize}
  \item \textbf{A} Each schedule word consumes four source bytes.
  \item \textbf{B} There are four compression rounds.
  \item \textbf{C} The byte count is stored in words.
\end{itemize}
